% Relatório do TCC2: edição mantida diretamente neste arquivo LaTeX.
\chapter{Materiais e métodos}\label{cap:tcc2-4}

Este capítulo descreve os recursos e procedimentos empregados no desenvolvimento. São apresentadas as ferramentas, a arquitetura de comunicação, as tecnologias de rede, a infraestrutura de observabilidade, a organização do projeto e os métodos de verificação e coleta. Os procedimentos de avaliação são apresentados neste capítulo, enquanto seus registros e resultados são reunidos na seção~\ref{sec:5.3}.

\section{Ferramentas e tecnologias}\label{sec:4.1}

A Godot 4.5 foi utilizada nos projetos Client e Server, com comportamentos programados em GDScript. O Blender integra o fluxo de autoria dos recursos visuais, e o projeto utiliza Git para controle de versão. Scripts PowerShell dão suporte à preparação dos diretórios compartilhados e à sincronização do mapa. Para a dashboard foram utilizados Python e FastAPI, com recepção UDP, comunicação WebSocket e consulta de dados do processo por psutil. A Tabela~\ref{tab:tcc2-1} resume o papel desses materiais.

\begingroup
\small
\setlength{\tabcolsep}{4pt}
\renewcommand{\arraystretch}{1.18}
\begin{longtable}{>{\raggedright\arraybackslash}p{\dimexpr0.32\textwidth-2\tabcolsep\relax}>{\raggedright\arraybackslash}p{\dimexpr0.68\textwidth-2\tabcolsep\relax}}
\caption{Ferramentas e responsabilidades}\label{tab:tcc2-1}\\
\toprule
Material & Utilização no trabalho \\
\midrule
\endfirsthead
\multicolumn{2}{l}{\small Ferramentas e responsabilidades (continuação)}\\
\toprule
Material & Utilização no trabalho \\
\midrule
\endhead
\midrule
\multicolumn{2}{r}{\small Continua na próxima página.}\\
\endfoot
\bottomrule
\endlastfoot
Godot e GDScript & Cenas, física, regras, interface, áudio e comunicação do jogo. \\
Blender & Autoria e manutenção dos modelos e recursos visuais. \\
Git & Organização e versionamento do projeto. \\
PowerShell & Preparação dos links compartilhados e geração de cenas do mapa de servidor. \\
Python e FastAPI & Recepção e disponibilização das amostras na dashboard. \\
psutil & Consulta de CPU e memória do processo identificado pelo PID. \\
\end{longtable}
\endgroup

Fonte: elaboração própria.

\section{Arquitetura e tecnologias de rede}\label{sec:4.2}

\subsection{Arquitetura de comunicação}\label{sec:4.2.1}

A arquitetura define como os participantes trocam informações e onde são decididos os resultados. Fiedler apresenta tanto a troca de comandos entre pares quanto a organização cliente-servidor. Nesta última, o servidor mantém uma referência comum para os clientes; no modelo de simulação determinística entre pares, a reprodução consistente dos comandos depende de condições iniciais e regras equivalentes. Esses modelos exigem soluções diferentes para sincronização. \cite{FIEDLER-REDE}

No projeto, adotou-se um processo servidor separado dos programas executados pelos jogadores. A entrada e a movimentação permanecem no Client, enquanto dano, recursos e pontuação são decididos pelo Server. Essa divisão combina resposta local ao controle com uma referência comum para os resultados das partidas. A organização dos componentes e o alcance dessa autoridade são detalhados na seção~\ref{sec:5.2}.

\subsection{Tecnologias de comunicação utilizadas}\label{sec:4.2.2}

O UDP é um protocolo de transporte baseado em datagramas. Sua especificação não garante entrega nem proteção contra duplicação de mensagens; requisitos adicionais devem ser tratados por outras camadas. No trabalho, ele aparece em dois contextos: como base da comunicação multijogador por ENet e como transporte direto das amostras enviadas à dashboard. A escolha do transporte, por si só, não determina se uma mensagem da aplicação será reenviada. \cite{UDP} \cite{GODOT-REDE}

A Godot disponibiliza comunicação de alto nível por chamadas remotas de procedimento, ou RPCs. Uma função pode ser executada em outro participante, com restrições sobre quem pode chamá-la e sobre a forma de transferência. Essa API permite usar a estrutura de nós do jogo para organizar a comunicação; o projeto a utiliza nos componentes que representam o servidor e cada lobby. \cite{GODOT-REDE}

Os modos de transferência usados são confiável, não confiável e não confiável ordenado. No primeiro, há confirmação e tentativas de reenvio; no segundo, uma mensagem pode se perder ou chegar fora de ordem; no terceiro, mensagens atrasadas podem ser descartadas quando uma atualização posterior já foi recebida. Essa diferença fundamenta o uso de eventos confiáveis para mudanças discretas e de atualizações substituíveis para movimento. \cite{GODOT-PEER}

Outra técnica utilizada foi a interpolação de snapshots. Um snapshot é um registro do estado em determinado instante. Ao manter um pequeno histórico dos registros recebidos, o Client pode apresentar uma posição intermediária entre dois deles, em vez de saltar diretamente para cada atualização. O buffer acrescenta um atraso de apresentação, em troca de uma transição visual mais contínua. \cite{FIEDLER-INTERPOLACAO}

\section{Observabilidade}\label{sec:4.3}

A observabilidade utiliza informações produzidas por um sistema para ajudar a compreender seu comportamento. A documentação do OpenTelemetry apresenta métricas como medições que descrevem aspectos de sua operação. No contexto deste jogo, elas permitem acompanhar condições como quantidade de participantes, tráfego de rede e tempo de execução de uma parte do servidor. O projeto possui instrumentação própria; não utiliza uma integração com o OpenTelemetry. \cite{OBSERVABILIDADE}

A infraestrutura combina apresentação em tempo real e gravação em CSV. A primeira permite acompanhar uma partida enquanto ocorre; a segunda preserva amostras para análise posterior. Para comparar condições de execução, o protocolo da seção~\ref{sec:4.8} prevê registrar o cenário, o intervalo de coleta e a origem dos indicadores.

\section{Processo de desenvolvimento}\label{sec:4.4}

O planejamento original do TCC1 adotou o Double Diamond como referência de processo. O modelo apresentado pelo Design Council organiza atividades de descoberta, definição, desenvolvimento e entrega. Neste trabalho, ele contextualiza o planejamento inicial, sem caracterizar um protocolo de avaliação das iterações realizadas. \cite{DOUBLE-DIAMOND}

O desenvolvimento integrou progressivamente os sistemas de jogador, combate, modos de partida, interface e observabilidade. O relatório apresenta essa organização por responsabilidades, relacionando as soluções adotadas aos problemas de interação e sincronização. A evolução do escopo é discutida na seção seguinte, e o funcionamento dos sistemas é detalhado no capítulo~\ref{cap:tcc2-5}.

\section{Evolução da proposta do TCC1}\label{sec:4.5}

A proposta original previa um participante hospedeiro e utilizava a denominação peer-to-peer para sua arquitetura. A implementação passou a separar o programa executado pelos jogadores do processo que mantém as partidas. Essa mudança permitiu executar o Server sem a apresentação do jogo e manter em um ponto comum as regras de dano, recursos, modos e resultados. RPCs e snapshots próprios substituíram a dependência do mecanismo de replicação automática descrito no protótipo inicial, tornando explícitas as responsabilidades de cada lado da comunicação.

A configuração livre por um jogador hospedeiro também foi substituída por presets compartilhados. Com isso, Client e Server trabalham com combinações previamente conhecidas de mapa, modo, tempo, pontuação e quantidade de participantes, e o gerenciador pode associar solicitações compatíveis ao mesmo lobby. Essa delimitação reduz as variações que precisam ser tratadas e verificadas na versão entregue.

O recorte de conteúdo foi ajustado para concentrar o desenvolvimento e a verificação em um conjunto menor de combinações. Foram disponibilizados TDM e KOTH em Highrise, enquanto Capture the Flag e o mapa Prototype ficaram fora da seleção principal. Highrise passou a reunir os elementos necessários aos dois modos entregues, como pontos de nascimento, recursos, jump pads e zona de controle. Dessa forma, o mesmo mapa pode ser utilizado no fluxo principal e nas duas modalidades disponíveis.

O catálogo final de combate também foi revisto. Pistola, SMG e shotgun assumem papéis diferenciados por cadência, dispersão e quantidade de traços, enquanto as granadas acrescentam dano em área. A melhoria de velocidade de movimentação prevista inicialmente deu lugar ao aumento temporário de cadência, integrado ao mesmo ciclo de armas, munição e feedback da interface. Essas alterações descrevem o produto efetivamente consolidado e são sintetizadas na seção~\ref{sec:5.3.3}.

\section{Organização do projeto}\label{sec:4.6}

O projeto foi organizado em dois projetos Godot independentes, Client e Server, e um diretório de código comum, \texttt{shared}. Essa estrutura permite executar o servidor separadamente da apresentação do jogo e reutilizar contratos de comunicação e definições nos dois lados. A Figura~\ref{fig:tcc2-2} apresenta os principais agrupamentos de responsabilidades, cujo funcionamento é discutido no relatório técnico.

O compartilhamento local é preparado por um script que cria links de diretório para uma mesma origem, evitando cópias independentes das definições comuns. Os arquivos de autoria dos assets são mantidos em pastas \texttt{import/}. A preparação desses recursos compõe o procedimento de geração dos executáveis descrito na seção~\ref{sec:5.2.11}.

Este relatório descreve a versão do projeto preparada para entrega. O código-fonte e a documentação técnica complementar estão disponíveis em \url{https://github.com/limaalan/TCC}. O Apêndice~\ref{ap:fontes-tecnicas} reúne os caminhos de consulta desse material.

\section{Método de verificação funcional}\label{sec:4.7}

A verificação funcional prevê executar o Server e conectar dois Clients, percorrendo o fluxo completo de partida. O roteiro inclui entrada no lobby, escolha de modo, nascimento dos personagens, combate, coleta de recursos, atualização da interface, encerramento e retorno ao menu. Para KOTH, é necessário observar separadamente a zona vazia, controlada e contestada, pois cada estado produz uma consequência diferente sobre a pontuação.

O registro do teste deve identificar a data, a plataforma, os executáveis utilizados e o resultado observado em cada cenário. A Tabela~\ref{tab:tcc2-5} organiza os comportamentos implementados e as situações que precisam ser exercitadas na versão preparada para entrega. Sua aprovação depende dos registros de execução, ainda assinalados como pendentes.

\begin{quote}
\small
\textbf{Evidência pendente:} registro do teste integrado da versão escolhida para a entrega, com data, identificação da build, participantes ou instâncias utilizadas e resultado de cada cenário. O procedimento acima está definido; sua aprovação não foi presumida.
\end{quote}

\section{Método de coleta e análise de métricas}\label{sec:4.8}

A infraestrutura de métricas permite coletar amostras periódicas, com intervalo padrão de um segundo, e gravá-las em CSV ou transmiti-las à dashboard. O coletor utiliza contadores do host ENet, estatísticas dos peers e informações de lobbies e jogadores. Os Clients também informam seus FPS, que são reunidos pelo Server. Essa infraestrutura fornece as variáveis para o experimento, mas a quantidade de amostras e a duração efetiva de cada rodada precisam ser registradas durante sua execução.

Para a avaliação desta versão, propõe-se comparar o servidor sem partidas ativas com uma partida de dois jogadores em TDM e uma de dois jogadores em KOTH. Cada rodada deve usar a mesma build, intervalo de coleta e ambiente, registrando o período de espera e o de jogo ativo. Experimentos com mais lobbies podem ampliar essa comparação até o limite configurado, desde que sejam identificados separadamente. Esse protocolo orienta a avaliação quantitativa e mantém as condições comparáveis entre os dois modos.

A análise deverá relacionar tráfego, RTT, perda estimada de pacotes e tempo medido do processamento dos lobbies ao número de participantes e ao estado da partida. Para evitar comparar medidas diferentes, devem ser preservados os nomes, as unidades e os timestamps do contrato. A medição de CPU e memória requer um registro adicional: esses valores são consultados pela dashboard no sistema operacional e não integram automaticamente os CSVs produzidos pelo logger do Godot.

\begin{quote}
\small
\textbf{Evidência pendente:} ambiente de execução, duração e repetições das rodadas, arquivos de dados brutos e resultados calculados. Não foram atribuídos valores a latência, CPU, memória ou capacidade do servidor sem essas medições.
\end{quote}
